% Summary based on the provided Deep Research PDF 
\section*{Kalman/State-Space Methods for Missing Data Imputation (State of the Art --- Summary)}

\subsection*{1. Existing tools, methods, and related work}
\textbf{Structural (state-space) time-series models} represent an observed series $y_t$ (e.g., 10-minute energy consumption) as a sum of interpretable latent components---level/trend, seasonality, and optional exogenous regressors (temperature, holidays, closure flags). In the linear-Gaussian case, inference is performed by the \textbf{Kalman filter} (forward recursion) and \textbf{Kalman smoother} (backward pass), yielding optimal (MMSE) state estimates and uncertainty quantification \cite{durbin2012,harvey1989}.

\textbf{Native missing-data handling.} In Kalman filtering, if $y_t$ is missing, the update step is skipped and the state is propagated using only the transition model; smoothing later uses surrounding observations to interpolate within gaps \cite{durbin2012}. This property makes SSMs a principled approach for irregular missingness and block outages.

\textbf{Seasonality representations.}
\begin{itemize}
\item \textit{Seasonal dummy states} (one state per intra-period position) provide a flexible seasonal shape but create a state dimension proportional to the period (e.g., 144 for daily, 1008 for weekly at 10-minute resolution), which can become computationally prohibitive \cite{durbin2012,harvey1989}.
\item \textit{Trigonometric/Fourier seasonal states} represent seasonality using a limited number of harmonics (sine/cosine pairs), dramatically reducing state dimension while preserving smooth periodic structure. This is central in TBATS for complex and multiple seasonalities \cite{delivera2011}.
\end{itemize}

\textbf{Practical software ecosystems.} Applied implementations commonly rely on optimized state-space toolkits (e.g., structural models and smoothers in state-space libraries; the PDF specifically references workflows akin to \texttt{statsmodels} structural components and R tooling such as \texttt{imputeTS} with Kalman smoothing) \cite{moritz2017}. These tools make Kalman-based imputation accessible, but their feasibility depends strongly on model dimensionality (especially seasonality design).

\subsection*{2. Critical analysis: strengths, limitations, and gaps}
\subsubsection*{Strengths}
\begin{itemize}
\item \textbf{Robustness to missingness patterns:} missing observations are handled without ad-hoc preprocessing by skipping measurement updates, while smoothing leverages past and future observations for gap reconstruction \cite{durbin2012}.
\item \textbf{Uncertainty estimates:} the smoother provides both point imputations and associated variances, enabling downstream systems to account for imputation confidence \cite{durbin2012}.
\item \textbf{Domain-informed structure:} explicit components (trend + multiple seasonalities + calendar/weather regressors) align with energy data drivers and support interpretable diagnostics \cite{harvey1989}.
\item \textbf{Competitive empirical performance in univariate settings:} applied studies and benchmarking reports often find Kalman-smoothing approaches strong for single-series time-series imputation compared to simpler baselines \cite{moritz2017,kannegowda2020}.
\end{itemize}

\subsubsection*{Limitations}
\begin{itemize}
\item \textbf{Model mis-specification risk:} imputations are only as good as the structural assumptions; regime changes (new equipment schedules, behavioral shifts) can bias reconstructions unless parameters and structure adapt over time \cite{harvey1989}.
\item \textbf{Long gaps and end-of-series (boundary) effects:} with no observations inside a long gap (or beyond the series endpoint), uncertainty grows and the method becomes effectively extrapolative; errors can propagate without correction until observations resume \cite{durbin2012}.
\item \textbf{Sensitivity to outliers under Gaussian noise:} spikes can distort state estimates; robust state-space variants exist but add complexity. For instance, robust mean-shift/outlier-adjusted formulations improve resilience to anomalies but can increase computational burden \cite{shankar2023}.
\item \textbf{Scalability and RAM usage for high-frequency multi-seasonality:} standard Kalman filtering scales poorly with state dimension ($O(m^2)$ memory and typically $O(m^3)$ time per step in naive forms). For 10-minute data with daily (144) and weekly (1008) dummy seasonals, $m$ can exceed $\sim$1100, leading to large covariance matrices and practical RAM blow-ups \cite{durbin2012}. This is a key deployment blocker in production-grade imputation at high resolution.
\end{itemize}

\subsubsection*{Key gaps and open needs for production pipelines}
\begin{itemize}
\item \textbf{Memory-efficient multi-seasonality at 10-minute resolution:} Fourier/trigonometric seasonality (as in TBATS) is a strong mitigation, but there is a practical need for principled harmonic selection and monitoring of seasonal drift under changing operational regimes \cite{delivera2011}.
\item \textbf{Approximate/structured filtering for large states:} low-rank or structured covariance approximations can reduce computational cost and may be necessary when multiple seasonalities and regressors still create large $m$ \cite{pnevmatikakis2014}.
\item \textbf{Robustness to anomaly+missingness interactions:} production systems need integrated strategies (outlier detection + robust filtering) to prevent anomalous points from contaminating imputations, while keeping compute bounded \cite{shankar2023}.
\item \textbf{Fallback strategies for extreme gaps and boundaries:} Kalman-based imputations should be complemented by guarded fallbacks (e.g., similarity-based historical substitution or simpler seasonal templates) when uncertainty becomes too large or gaps occur at the series end (where smoothing cannot use future data) \cite{durbin2012}.
\end{itemize}

\subsection*{3. Condensed comparison of main model variants}
\begin{table}[h]
\centering
\small
\begin{tabular}{p{3.1cm}p{2.7cm}p{2.8cm}p{4.2cm}}
\hline
\textbf{Variant} & \textbf{Missing data} & \textbf{Seasonality} & \textbf{Scalability / notes} \
\hline
Structural SSM + dummy seasonal & Native (skip update, smooth) \cite{durbin2012} & One state per intra-period & Often infeasible for 10-min + weekly due to large $m$ and covariance RAM \cite{durbin2012}. \
Trigonometric seasonal SSM & Native \cite{durbin2012} & Fourier harmonics & Much smaller $m$; practical for daily+weekly patterns \cite{delivera2011}. \
TBATS-style multi-seasonality & Native \cite{delivera2011} & Multiple Fourier seasonalities & Good for complex seasonal structure; complexity rises with many harmonics \cite{delivera2011}. \
Robust/outlier-adjusted SSM & Native \cite{shankar2023} & Any (dummy/Fourier) & Better anomaly resilience; added tuning/compute, sometimes higher dimensionality \cite{shankar2023}. \
\hline
\end{tabular}
\end{table}

\subsection*{4. Takeaway for high-frequency energy imputation}
Kalman/SSM imputation is state-of-the-art when the model matches the data-generating mechanisms (trend + calendar/weather effects + multiple seasonalities), and it is particularly attractive because missingness is handled natively and uncertainty is explicit \cite{durbin2012,harvey1989}. The primary practical obstacle for 10-minute energy data is \textbf{state explosion} from naive seasonal encodings, motivating Fourier seasonality (TBATS-style) and, when needed, approximate/structured filtering to keep RAM and compute within operational constraints \cite{delivera2011,pnevmatikakis2014}.

\begin{thebibliography}{99}
\bibitem{durbin2012}
J.~Durbin and S.~J.~Koopman, \textit{Time Series Analysis by State Space Methods}, 2nd~ed. Oxford University Press, 2012.

\bibitem{harvey1989}
A.~C.~Harvey, \textit{Forecasting, Structural Time Series Models and the Kalman Filter}. Cambridge University Press, 1989.

\bibitem{delivera2011}
C.~M.~De~Livera, R.~J.~Hyndman, and R.~D.~Snyder, ``Forecasting time series with complex seasonal patterns using exponential smoothing: the TBATS approach,'' \textit{Journal of the American Statistical Association}, vol.~106, no.~496, pp.~1513--1527, 2011, doi:10.1198/jasa.2011.ap10271.

\bibitem{moritz2017}
S.~Moritz and T.~Bartz-Beielstein, ``imputeTS: Time series missing value imputation in R,'' \textit{The R Journal}, vol.~9, no.~1, pp.~207--218, 2017, doi:10.32614/RJ-2017-009.

\bibitem{pnevmatikakis2014}
E.~A.~Pnevmatikakis, K.~R.~Rad, J.~Huggins, and L.~Paninski, ``Fast Kalman filtering and forward-backward smoothing via a low-rank perturbative approach,'' \textit{Journal of Computational and Graphical Statistics}, vol.~23, no.~2, pp.~316--338, 2014, doi:10.1080/10618600.2012.760461.

\bibitem{kannegowda2020}
B.~Kannegowda \textit{et al.}, ``Comparative analysis of univariate and multivariate methods for imputing missing precipitation data in a tropical region,'' \textit{Environmental Progress & Sustainable Energy}, vol.~39, no.~4, 2020, doi:10.1002/ep.13451.

\bibitem{shankar2023}
R.~Shankar \textit{et al.}, ``Robust outlier-adjusted mean-shift estimation of state-space models,'' arXiv:2511.15155, 2023.
\end{thebibliography}
